\section{Bounded-weight shallow teachers from Fourier coefficients}
\label{sec:teacher}

The target used above has a small teacher at the same shallow depth as
 the learner.  We prove this using a general construction that may also be
 useful for other targets.  All networks in this section have one tanh
 hidden layer, one constant input coordinate, and a linear readout.

For $S\subseteq\{1,\ldots,d\}$, write
$\chi_S(x)=\prod_{i\in S}x_i$.  A Boolean function has the Fourier expansion
\[
 f(x)=\sum_{S\subseteq\{1,\ldots,d\}}\widehat f(S)\chi_S(x),
 \qquad
 \mathcal A(f)=\sum_S|\widehat f(S)|.
\]
The Fourier coefficients here use the uniform measure on the Boolean
cube; the learning theorem does not require a uniform input distribution.

\begin{theorem}[A bounded-weight shallow teacher]\label{thm:fourier-teacher}
Every $f:\{-1,+1\}^d\to\{-1,+1\}$ has a one-hidden-layer tanh
implementation that uniformly approximates $f$ within $1/4$, has every
weight of absolute value at most $1$, and has hidden width at most
\begin{equation}\label{eq:teacher-width}
 \left\lceil
 288\mathcal A(f)^2(2d+1)(d+2)\log 2
 \right\rceil.
\end{equation}
In particular, its signed output is greater than $3/4$ at every input.
\end{theorem}

\begin{lemma}[Parity correlation of bounded tanh features]
\label{lem:parity-feature}
For $r\ge0$, let $b_r=0$ when $r$ is odd and $b_r=1$ when $r$ is even.
For independent uniform signs $s_1,\ldots,s_r$, define
\[
 c_r=\mathbb E_s\left[
    \left(\prod_{i=1}^r s_i\right)
    \tanh\left(\sum_{i=1}^r s_i+b_r\right)
 \right].
\]
Then
\begin{equation}\label{eq:parity-correlation}
 |c_r|\ge\frac{1}{3\sqrt{2r+1}}.
\end{equation}
The empty product and sum give $c_0=\tanh 1$.
\end{lemma}
\begin{proof}
Put $m=\lfloor r/2\rfloor$ and
$\beta_m=\binom{2m}{m}/4^m$.  The highest-degree Fourier coefficient of
$\operatorname{sign}(\sum_i s_i+b_r)$ is $(-1)^m\beta_m$.
For $r\ge1$, this follows by splitting the binomial sum at
$\lfloor r/2\rfloor$ and using
\[
 \sum_{j=0}^{q}(-1)^j\binom rj
   =(-1)^q\binom{r-1}{q};
\]
the case $r=0$ is immediate.

Every possible value of $T=\sum_i s_i+b_r$ is an odd integer, and each
individual value has probability at most $\beta_m$.  For even $r$ this
is the central-binomial bound.  For odd $r=2m+1$, the largest probability
is $\frac{2m+1}{2m+2}\beta_m$.  Since
$|\operatorname{sign}(t)-\tanh t|\le2e^{-2|t|}$ for $t\ne0$,
\begin{align*}
 \left|c_r-(-1)^m\beta_m\right|
 &\le\mathbb E\left|
     \operatorname{sign}(T)-\tanh T\right|\\
 &\le4\beta_m\sum_{j\ge0}e^{-2(2j+1)}
   =\frac{4e^{-2}}{1-e^{-4}}\beta_m
   <\frac{7}{12}\beta_m.
\end{align*}
The last inequality follows from $e^2>7$ and
$4e^{-2}/(1-e^{-4})=4e^2/(e^4-1)<7/12$.
Thus $|c_r|>\frac{5}{12}\beta_m>\frac13\beta_m$.

Finally, $\beta_m\ge1/\sqrt{4m+1}$.  This follows by induction from
$\beta_0=1$ and
$\beta_{m+1}/\beta_m=(2m+1)/(2m+2)$: after squaring, the induction step is
\[
 (2m+1)^2(4m+5)\ge4(m+1)^2(4m+1),
\]
whose left-hand side exceeds the right-hand side by $1$.
Since $4m+1\le2r+1$, the claimed bound follows.
\end{proof}

\begin{proof}[Proof of Theorem~\ref{thm:fourier-teacher}]
For a subset $S$ of cardinality $r$ and a sign vector $s$ indexed by $S$,
let
\[
 h_{S,s}(x)=\tanh\left(\sum_{i\in S}s_i x_i+b_r\right).
\]
Every affine coefficient belongs to $\{0,-1,+1\}$, including the
coefficient of the constant input.  Changing variables from $s_i$ to
$s_i x_i$ gives the exact identity
\begin{equation}\label{eq:parity-mixture}
 \mathbb E_s\left[
   \left(\prod_{i\in S}s_i\right)h_{S,s}(x)
 \right]=c_r\chi_S(x).
\end{equation}
Consequently, the Fourier expansion expresses $f$ as
\[
 f(x)=B\mathbb E[H(x)],
 \qquad
 B=\sum_S\frac{|\widehat f(S)|}{|c_{|S|}|}
   \le3\mathcal A(f)\sqrt{2d+1},
\]
where $H$ is a signed feature of the form $h_{S,s}$ and $|H(x)|\le1$.
Explicitly, choose $S$ with probability proportional to
$|\widehat f(S)|/|c_{|S|}|$, choose $s$ uniformly, and multiply
$h_{S,s}$ by the signs prescribed by
$\widehat f(S)/c_{|S|}$ and $\prod_{i\in S}s_i$.
The empty subset uses the constant feature $\tanh1$ and needs no separate
construction.

Take
\[
 n=\left\lceil32B^2\log(4\cdot2^d)\right\rceil
\]
independent features $H_1,\ldots,H_n$.  Hoeffding's inequality and a union
bound over the cube give
\[
 \mathbb P\left\{
  \max_x\left|\frac Bn\sum_{j=1}^nH_j(x)-f(x)\right|\ge\frac14
 \right\}
 \le2\cdot2^d\exp\left(-\frac{n}{32B^2}\right)
 \le\frac12.
\]
A choice of features therefore achieves uniform error strictly below
$1/4$.  Put each feature's sign in its readout coefficient, whose
magnitude is $B/n$.  Because $\mathcal A(f)\ge1$ and $|c_r|\le1$, we have
$B\ge1$ and $n\ge B$; hence every readout weight also has magnitude at
most $1$.  Finally,
$B^2\le9\mathcal A(f)^2(2d+1)$ and
$\log(4\cdot2^d)=(d+2)\log2$ give~\eqref{eq:teacher-width}.
\end{proof}

\begin{corollary}[A shallow teacher for the selector target]
\label{cor:selector-teacher}
For the target $f_k$ on $d=2k+1$ inputs, equal to $+1$ precisely when
$z=-1$ and $X\ge Y$, there is a one-hidden-layer tanh teacher with all
weights bounded by $1$, global signed output greater than $3/4$, and
width at most
\begin{equation}\label{eq:selector-teacher-width}
 W_k=\left\lceil
 288(k+2)^2(4k+3)(2k+3)\log2
 \right\rceil=O(k^4).
\end{equation}
In the convention that $C_2$ bounds both hidden width and individual
weight magnitudes, $C_2(f_k)\le W_k$.
\end{corollary}
\begin{proof}
Let $x_i,y_i$ be the two integers' Boolean bits, with $i=0$ the least
significant position, and put $E_i=(1+x_i y_i)/2$.  The comparison sign
$G$, equal to $+1$ when $X\ge Y$ and $-1$ otherwise, satisfies
\[
 G(x,y)=\prod_{i=0}^{k-1}E_i
   +\sum_{i=0}^{k-1}\frac{x_i-y_i}{2}
       \prod_{j=i+1}^{k-1}E_j.
\]
The first term handles equality; otherwise the highest differing bit
contributes the comparison sign and all other summands vanish.
Each $E_i$ and $(x_i-y_i)/2$ has Fourier $\ell^1$ norm $1$.
Submultiplicativity and the triangle inequality therefore give
$\mathcal A(G)\le k+1$.  The selector target is
\[
 f_k(x,y,z)=\frac{1-z}{2}G(x,y)-\frac{1+z}{2},
\]
so $\mathcal A(f_k)\le k+2$.  Apply
Theorem~\ref{thm:fourier-teacher} with $d=2k+1$.
\end{proof}
